% alias=汉语字体设置
%%%%%
% 此文件提供了“！！本地！！”的文件中文字体部署。若整个文档中不含有汉语文字则不需要包含此宏包
% \ProvidesPackage{/Users/Disk/03-Com-Scientia/01-TeX/Settings/Fonts/SimplifiedChineseFonts}
\usepackage[fontset=none, scheme=plain]{ctex}
%\usepackage[fontset=none]{ctex}
% directory of fonts: /Users/veritas/Library/Fonts


%%%%% <main font>
% 主字体：{Founder:方正新书宋_GBK	Regular}
%   标识符：FZXSSK--GBK1-0
%   文件：FZXSSK.TTF
% boldfont: {Founder:方正粗宋简繁}
%   标识符：FZCSK--GBK1-0
%   文件：FZCSK.TTF
% italicfont: {Founder:方正汉文正楷 简}
%   标识符：FZHANWZKS--GB1-0
%   文件：FZHanWZKJW.TTF

\setCJKmainfont[BoldFont={FZCSK--GBK1-0}, ItalicFont={FZHANWZKS--GB1-0},]{FZXSSK--GBK1-0}
%%%%% </main font>
% Checked---------------------%
%%%%% <sans serif font>


% 方正黑体_GBK
\setCJKsansfont[BoldFont={FZCHSJW--GB1-0}]{FZHTK--GBK1-0}

%%%%% </sans serif font>
% <font-families>
% 仿宋字体：{Founder:方正仿宋_GBK	Regular}
%   标识符：FZFSK--GBK1-0
%   文件：FZFSK.TTF
\setCJKmonofont{FZFSK--GBK1-0}
\setCJKfamilyfont{zhfangsong}{FZFSJW--GB1-0}
\newcommand*{\fangsong}{\CJKfamily{zhfangsong}}




% </font-families>
\ctexset{space=auto} % Automatic chinese & english word spacing